% --- Archon physics formalization source begin ---
% archon:physics
% archon:covers PhyXMiniProblems/problem_phyx_mini_0752.lean
% archon:source-report reports/phyx_mini/problem_phyx_mini_0752.source.json

\chapter{Physics problem phyx\_mini\_0752}
\label{ch:PhyXMiniProblems_problem_phyx_mini_0752}

\paragraph{Problem source.}
\#\# Physical scenario

During volcanic eruptions, chunks of solid rock can be blasted out of the volcano; these projectiles are called volcanic bombs. Figure shows a cross section of Mt. Fuji, in Japan.

\#\# Question

At what initial speed would a bomb have to be ejected, at angle \$\textbackslash{}theta\_0 = 35\textasciicircum{}\textbackslash{}circ\$ to the horizontal, from the vent at \$A\$ in order to fall at the foot of the volcano at \$B\$, at vertical distance \$h = 3.30\textbackslash{} km\$ and horizontal distance \$d = 9.40\textbackslash{} km\$? Ignore, for the moment, the effect of air on the bomb's travel.

\#\# Answer choices

- A: A: 235.5 m/s

- B: B: 245.5 m/s

- C: C: 255.5 m/s

- D: D: 265.5 m/s

\#\# Figure caption (auxiliary; use the image as primary evidence)

The image depicts a projectile motion scenario involving a volcano.

- There is a volcano with a vent labeled "A" at the top. From this point, a projectile is launched.
- The projectile follows a curved trajectory towards a point labeled "B" on the ground.
- The initial angle of launch is labeled \textbackslash{}(\textbackslash{}theta\_0\textbackslash{}) and is marked with a double-headed arrow.
- The height of the volcano is labeled "h" with an arrow indicating the vertical distance from the base to the top.
- The horizontal distance from point "A" to point "B" is labeled "d".
- Red streaks represent the explosive force at the vent.
- The path of the projectile is indicated by a curved brown line.

This diagram illustrates the mechanics of a projectile launched from a raised elevation, showing key parameters like launch angle, height, and range.

\#\# Dataset metadata

Kinematics; Physical Model Grounding Reasoning, Multi-Formula Reasoning

\paragraph{Recorded answer/context.}
C: C: 255.5 m/s

\paragraph{Figure/image path.}
/root/proposal\_for\_physic/science-mango/phyx\_mini\_run/phyx\_data/test\_image/752.png

\paragraph{Formalization target.}
create a compiling Lean file with sorry bodies at `PhyXMiniProblems/problem\_phyx\_mini\_0752.lean`. The Lean declarations must preserve the physical quantities, dimensions or dimensional roles, figure labels, governing-law hypotheses, and final relation expressed by this problem.
Use Mathlib/Physlib names found through LeanExplore where available. If a physics API is missing, introduce faithful local abstractions rather than scalar placeholder aliases.

\begin{theorem}[Physics formalization target]
\label{thm:physics:phyx_mini_0752:target}
\lean{PhyXMiniProblems.ProblemPhyXMini0752.problem_phyx_mini_0752}
\uses{def:physics:phyx-mini-0752:phyxminiproblems-problemphyxmini0752-speedinmeterspersecond, def:physics:phyx-mini-0752:phyxminiproblems-problemphyxmini0752-volcanicbombsetup, def:physics:phyx-mini-0752:phyxminiproblems-problemphyxmini0752-matchesprimaryfigure, def:physics:phyx-mini-0752:phyxminiproblems-problemphyxmini0752-matchesproblemdescription, def:physics:phyx-mini-0752:phyxminiproblems-problemphyxmini0752-usesstandardnearearthgravity, def:physics:phyx-mini-0752:phyxminiproblems-problemphyxmini0752-hasphysicalprojectileparameters, def:physics:phyx-mini-0752:phyxminiproblems-problemphyxmini0752-satisfiesuniformgravityprojectilelaws, def:physics:phyx-mini-0752:phyxminiproblems-problemphyxmini0752-requiredlaunchspeedinmeterspersecond, def:physics:phyx-mini-0752:phyxminiproblems-problemphyxmini0752-matchesdisplayedlaunchspeed, def:physics:phyx-mini-0752:phyxminiproblems-problemphyxmini0752-isuniqueclosestlaunchspeedchoice}
The assigned autoformalize agent should translate this physics problem into Lean declarations in the covered file, with theorem and lemma proof bodies written as `by sorry`.
\end{theorem}
\begin{proof}
This is an autoformalization task, not a proof task. Produce statements that can later be proved by the physics prover without weakening the physical model.
\end{proof}

% --- Archon declaration topology begin ---
\section{Declaration topology}
\begin{definition}[acceleration Dimension]
  \label{def:physics:phyx-mini-0752:phyxminiproblems-problemphyxmini0752-accelerationdimension}
  \lean{PhyXMiniProblems.ProblemPhyXMini0752.accelerationDimension}
  The physical dimension of acceleration, L T\textasciicircum{}-2
\end{definition}
\begin{proof}
  This is the declared model definition of acceleration Dimension.
\end{proof}
\begin{definition}[Length Quantity]
  \label{def:physics:phyx-mini-0752:phyxminiproblems-problemphyxmini0752-lengthquantity}
  \lean{PhyXMiniProblems.ProblemPhyXMini0752.LengthQuantity}
  A nonnegative physical length
\end{definition}
\begin{proof}
  This is the declared model definition of Length Quantity.
\end{proof}
\begin{definition}[Signed Length Quantity]
  \label{def:physics:phyx-mini-0752:phyxminiproblems-problemphyxmini0752-signedlengthquantity}
  \lean{PhyXMiniProblems.ProblemPhyXMini0752.SignedLengthQuantity}
  A signed one-dimensional physical displacement
\end{definition}
\begin{proof}
  This is the declared model definition of Signed Length Quantity.
\end{proof}
\begin{definition}[Duration Quantity]
  \label{def:physics:phyx-mini-0752:phyxminiproblems-problemphyxmini0752-durationquantity}
  \lean{PhyXMiniProblems.ProblemPhyXMini0752.DurationQuantity}
  A nonnegative duration elapsed since launch
\end{definition}
\begin{proof}
  This is the declared model definition of Duration Quantity.
\end{proof}
\begin{definition}[Speed Quantity]
  \label{def:physics:phyx-mini-0752:phyxminiproblems-problemphyxmini0752-speedquantity}
  \lean{PhyXMiniProblems.ProblemPhyXMini0752.SpeedQuantity}
  A nonnegative physical launch-speed magnitude
\end{definition}
\begin{proof}
  This is the declared model definition of Speed Quantity.
\end{proof}
\begin{definition}[Acceleration Quantity]
  \label{def:physics:phyx-mini-0752:phyxminiproblems-problemphyxmini0752-accelerationquantity}
  \lean{PhyXMiniProblems.ProblemPhyXMini0752.AccelerationQuantity}
  \uses{def:physics:phyx-mini-0752:phyxminiproblems-problemphyxmini0752-accelerationdimension}
  A nonnegative magnitude of gravitational acceleration
\end{definition}
\begin{proof}
  This is the declared model definition of Acceleration Quantity.
\end{proof}
\begin{definition}[length In Meters]
  \label{def:physics:phyx-mini-0752:phyxminiproblems-problemphyxmini0752-lengthinmeters}
  \lean{PhyXMiniProblems.ProblemPhyXMini0752.lengthInMeters}
  \uses{def:physics:phyx-mini-0752:phyxminiproblems-problemphyxmini0752-lengthquantity}
  Metre readout of a nonnegative physical length
\end{definition}
\begin{proof}
  This is the declared model definition of length In Meters.
\end{proof}
\begin{definition}[signed Length In Meters]
  \label{def:physics:phyx-mini-0752:phyxminiproblems-problemphyxmini0752-signedlengthinmeters}
  \lean{PhyXMiniProblems.ProblemPhyXMini0752.signedLengthInMeters}
  \uses{def:physics:phyx-mini-0752:phyxminiproblems-problemphyxmini0752-signedlengthquantity}
  Metre readout of a signed physical displacement
\end{definition}
\begin{proof}
  This is the declared model definition of signed Length In Meters.
\end{proof}
\begin{definition}[duration In Seconds]
  \label{def:physics:phyx-mini-0752:phyxminiproblems-problemphyxmini0752-durationinseconds}
  \lean{PhyXMiniProblems.ProblemPhyXMini0752.durationInSeconds}
  \uses{def:physics:phyx-mini-0752:phyxminiproblems-problemphyxmini0752-durationquantity}
  Second readout of a physical duration
\end{definition}
\begin{proof}
  This is the declared model definition of duration In Seconds.
\end{proof}
\begin{definition}[speed In Meters Per Second]
  \label{def:physics:phyx-mini-0752:phyxminiproblems-problemphyxmini0752-speedinmeterspersecond}
  \lean{PhyXMiniProblems.ProblemPhyXMini0752.speedInMetersPerSecond}
  \uses{def:physics:phyx-mini-0752:phyxminiproblems-problemphyxmini0752-speedquantity}
  Metre-per-second readout of a physical speed
\end{definition}
\begin{proof}
  This is the declared model definition of speed In Meters Per Second.
\end{proof}
\begin{definition}[acceleration In Meters Per Second Squared]
  \label{def:physics:phyx-mini-0752:phyxminiproblems-problemphyxmini0752-accelerationinmeterspersecondsquared}
  \lean{PhyXMiniProblems.ProblemPhyXMini0752.accelerationInMetersPerSecondSquared}
  \uses{def:physics:phyx-mini-0752:phyxminiproblems-problemphyxmini0752-accelerationquantity}
  Metre-per-second-squared readout of an acceleration magnitude
\end{definition}
\begin{proof}
  This is the declared model definition of acceleration In Meters Per Second Squared.
\end{proof}
\begin{definition}[Projectile Kind]
  \label{def:physics:phyx-mini-0752:phyxminiproblems-problemphyxmini0752-projectilekind}
  \lean{PhyXMiniProblems.ProblemPhyXMini0752.ProjectileKind}
  The physical object launched by the eruption
\end{definition}
\begin{proof}
  This is the declared model definition of Projectile Kind.
\end{proof}
\begin{definition}[Projectile Model]
  \label{def:physics:phyx-mini-0752:phyxminiproblems-problemphyxmini0752-projectilemodel}
  \lean{PhyXMiniProblems.ProblemPhyXMini0752.ProjectileModel}
  The approximation explicitly requested in the problem
\end{definition}
\begin{proof}
  This is the declared model definition of Projectile Model.
\end{proof}
\begin{definition}[Physical Location]
  \label{def:physics:phyx-mini-0752:phyxminiproblems-problemphyxmini0752-physicallocation}
  \lean{PhyXMiniProblems.ProblemPhyXMini0752.PhysicalLocation}
  Named physical locations in the prose and primary image
\end{definition}
\begin{proof}
  This is the declared model definition of Physical Location.
\end{proof}
\begin{definition}[Figure Point]
  \label{def:physics:phyx-mini-0752:phyxminiproblems-problemphyxmini0752-figurepoint}
  \lean{PhyXMiniProblems.ProblemPhyXMini0752.FigurePoint}
  Distinguished points needed to transcribe the diagram's geometry
\end{definition}
\begin{proof}
  This is the declared model definition of Figure Point.
\end{proof}
\begin{definition}[Figure Object]
  \label{def:physics:phyx-mini-0752:phyxminiproblems-problemphyxmini0752-figureobject}
  \lean{PhyXMiniProblems.ProblemPhyXMini0752.FigureObject}
  Physical or graphical objects visible in image 752.png
\end{definition}
\begin{proof}
  This is the declared model definition of Figure Object.
\end{proof}
\begin{definition}[Figure Label]
  \label{def:physics:phyx-mini-0752:phyxminiproblems-problemphyxmini0752-figurelabel}
  \lean{PhyXMiniProblems.ProblemPhyXMini0752.FigureLabel}
  Text and mathematical labels printed in the primary image
\end{definition}
\begin{proof}
  This is the declared model definition of Figure Label.
\end{proof}
\begin{definition}[Mt Fuji Projectile Figure]
  \label{def:physics:phyx-mini-0752:phyxminiproblems-problemphyxmini0752-mtfujiprojectilefigure}
  \lean{PhyXMiniProblems.ProblemPhyXMini0752.MtFujiProjectileFigure}
  \uses{def:physics:phyx-mini-0752:phyxminiproblems-problemphyxmini0752-figurepoint, def:physics:phyx-mini-0752:phyxminiproblems-problemphyxmini0752-figureobject, def:physics:phyx-mini-0752:phyxminiproblems-problemphyxmini0752-figurelabel}
  Literal qualitative and schematic-coordinate information from the raster. The coordinates are dimensionless drawing coordinates, not physical lengths
\end{definition}
\begin{proof}
  This is the declared model definition of Mt Fuji Projectile Figure.
\end{proof}
\begin{definition}[Volcanic Bomb Setup]
  \label{def:physics:phyx-mini-0752:phyxminiproblems-problemphyxmini0752-volcanicbombsetup}
  \lean{PhyXMiniProblems.ProblemPhyXMini0752.VolcanicBombSetup}
  \uses{def:physics:phyx-mini-0752:phyxminiproblems-problemphyxmini0752-lengthquantity, def:physics:phyx-mini-0752:phyxminiproblems-problemphyxmini0752-signedlengthquantity, def:physics:phyx-mini-0752:phyxminiproblems-problemphyxmini0752-durationquantity, def:physics:phyx-mini-0752:phyxminiproblems-problemphyxmini0752-speedquantity, def:physics:phyx-mini-0752:phyxminiproblems-problemphyxmini0752-accelerationquantity, def:physics:phyx-mini-0752:phyxminiproblems-problemphyxmini0752-projectilekind, def:physics:phyx-mini-0752:phyxminiproblems-problemphyxmini0752-projectilemodel, def:physics:phyx-mini-0752:phyxminiproblems-problemphyxmini0752-physicallocation, def:physics:phyx-mini-0752:phyxminiproblems-problemphyxmini0752-mtfujiprojectilefigure}
  Independent physical quantities for the launch. The displacement functions use an origin at vent A, with positive horizontal direction toward B and positive vertical direction upward. In particular, launchSpeed is not defined from an answer choice or from the closed-form result
\end{definition}
\begin{proof}
  This is the declared model definition of Volcanic Bomb Setup.
\end{proof}
\begin{definition}[Matches Primary Figure]
  \label{def:physics:phyx-mini-0752:phyxminiproblems-problemphyxmini0752-matchesprimaryfigure}
  \lean{PhyXMiniProblems.ProblemPhyXMini0752.MatchesPrimaryFigure}
  \uses{def:physics:phyx-mini-0752:phyxminiproblems-problemphyxmini0752-volcanicbombsetup}
  Primary-image evidence: A lies vertically above the left endpoint of the horizontal d arrow, B is at the same level as that endpoint and to its right, the curve runs from A to B, and the launch arrow points upward and right at the angle marked theta\_0
\end{definition}
\begin{proof}
  This is the declared model definition of Matches Primary Figure.
\end{proof}
\begin{definition}[degrees]
  \label{def:physics:phyx-mini-0752:phyxminiproblems-problemphyxmini0752-degrees}
  \lean{PhyXMiniProblems.ProblemPhyXMini0752.degrees}
  Convert a degree readout to Mathlib's physical angle modulo 2 * pi
\end{definition}
\begin{proof}
  This is the declared model definition of degrees.
\end{proof}
\begin{definition}[Matches Problem Description]
  \label{def:physics:phyx-mini-0752:phyxminiproblems-problemphyxmini0752-matchesproblemdescription}
  \lean{PhyXMiniProblems.ProblemPhyXMini0752.MatchesProblemDescription}
  \uses{def:physics:phyx-mini-0752:phyxminiproblems-problemphyxmini0752-lengthinmeters, def:physics:phyx-mini-0752:phyxminiproblems-problemphyxmini0752-volcanicbombsetup, def:physics:phyx-mini-0752:phyxminiproblems-problemphyxmini0752-degrees}
  The qualitative scenario and numerical readouts stated in the question. The kilometre data are converted to coherent-SI metre readouts. No launch-speed value occurs here
\end{definition}
\begin{proof}
  This is the declared model definition of Matches Problem Description.
\end{proof}
\begin{definition}[Uses Standard Near Earth Gravity]
  \label{def:physics:phyx-mini-0752:phyxminiproblems-problemphyxmini0752-usesstandardnearearthgravity}
  \lean{PhyXMiniProblems.ProblemPhyXMini0752.UsesStandardNearEarthGravity}
  \uses{def:physics:phyx-mini-0752:phyxminiproblems-problemphyxmini0752-accelerationinmeterspersecondsquared, def:physics:phyx-mini-0752:phyxminiproblems-problemphyxmini0752-volcanicbombsetup}
  The standard near-Earth gravitational acceleration used by the textbook projectile model. This calibration does not constrain the answer speed
\end{definition}
\begin{proof}
  This is the declared model definition of Uses Standard Near Earth Gravity.
\end{proof}
\begin{definition}[Has Physical Projectile Parameters]
  \label{def:physics:phyx-mini-0752:phyxminiproblems-problemphyxmini0752-hasphysicalprojectileparameters}
  \lean{PhyXMiniProblems.ProblemPhyXMini0752.HasPhysicalProjectileParameters}
  \uses{def:physics:phyx-mini-0752:phyxminiproblems-problemphyxmini0752-lengthinmeters, def:physics:phyx-mini-0752:phyxminiproblems-problemphyxmini0752-durationinseconds, def:physics:phyx-mini-0752:phyxminiproblems-problemphyxmini0752-speedinmeterspersecond, def:physics:phyx-mini-0752:phyxminiproblems-problemphyxmini0752-accelerationinmeterspersecondsquared, def:physics:phyx-mini-0752:phyxminiproblems-problemphyxmini0752-volcanicbombsetup}
  Positivity and direction assumptions selecting the physical upward/rightward branch shown in the figure. None assigns a numerical launch speed
\end{definition}
\begin{proof}
  This is the declared model definition of Has Physical Projectile Parameters.
\end{proof}
\begin{definition}[Satisfies Uniform Gravity Projectile Laws]
  \label{def:physics:phyx-mini-0752:phyxminiproblems-problemphyxmini0752-satisfiesuniformgravityprojectilelaws}
  \lean{PhyXMiniProblems.ProblemPhyXMini0752.SatisfiesUniformGravityProjectileLaws}
  \uses{def:physics:phyx-mini-0752:phyxminiproblems-problemphyxmini0752-durationquantity, def:physics:phyx-mini-0752:phyxminiproblems-problemphyxmini0752-lengthinmeters, def:physics:phyx-mini-0752:phyxminiproblems-problemphyxmini0752-signedlengthinmeters, def:physics:phyx-mini-0752:phyxminiproblems-problemphyxmini0752-durationinseconds, def:physics:phyx-mini-0752:phyxminiproblems-problemphyxmini0752-speedinmeterspersecond, def:physics:phyx-mini-0752:phyxminiproblems-problemphyxmini0752-accelerationinmeterspersecondsquared, def:physics:phyx-mini-0752:phyxminiproblems-problemphyxmini0752-volcanicbombsetup}
  The standard no-drag projectile equations in coherent SI readouts. Horizontal motion is uniform; vertical motion has constant downward acceleration. At the positive flight duration, the bomb's displacement is d horizontally and -h vertically, exactly encoding landing at B below A. These are governing laws and endpoint geometry only. They contain neither a solved launch-speed formula nor any displayed answer value
\end{definition}
\begin{proof}
  This is the declared model definition of Satisfies Uniform Gravity Projectile Laws.
\end{proof}
\begin{definition}[required Launch Speed In Meters Per Second]
  \label{def:physics:phyx-mini-0752:phyxminiproblems-problemphyxmini0752-requiredlaunchspeedinmeterspersecond}
  \lean{PhyXMiniProblems.ProblemPhyXMini0752.requiredLaunchSpeedInMetersPerSecond}
  \uses{def:physics:phyx-mini-0752:phyxminiproblems-problemphyxmini0752-lengthinmeters, def:physics:phyx-mini-0752:phyxminiproblems-problemphyxmini0752-accelerationinmeterspersecondsquared, def:physics:phyx-mini-0752:phyxminiproblems-problemphyxmini0752-volcanicbombsetup}
  Eliminating flight time from the two endpoint equations gives this unrounded speed readout. It depends only on the independent distance, angle, and gravity fields, not on setup.launchSpeed and not on an answer choice
\end{definition}
\begin{proof}
  This is the declared model definition of required Launch Speed In Meters Per Second.
\end{proof}
\begin{definition}[Answer Choice]
  \label{def:physics:phyx-mini-0752:phyxminiproblems-problemphyxmini0752-answerchoice}
  \lean{PhyXMiniProblems.ProblemPhyXMini0752.AnswerChoice}
  Labels of the four speed choices printed with the problem
\end{definition}
\begin{proof}
  This is the declared model definition of Answer Choice.
\end{proof}
\begin{definition}[speed Meters Per Second]
  \label{def:physics:phyx-mini-0752:phyxminiproblems-problemphyxmini0752-answerchoice-speedmeterspersecond}
  \lean{PhyXMiniProblems.ProblemPhyXMini0752.AnswerChoice.speedMetersPerSecond}
  \uses{def:physics:phyx-mini-0752:phyxminiproblems-problemphyxmini0752-answerchoice}
  Metres-per-second value printed beside each answer choice
\end{definition}
\begin{proof}
  This is the declared model definition of speed Meters Per Second.
\end{proof}
\begin{definition}[recorded Answer Choice]
  \label{def:physics:phyx-mini-0752:phyxminiproblems-problemphyxmini0752-recordedanswerchoice}
  \lean{PhyXMiniProblems.ProblemPhyXMini0752.recordedAnswerChoice}
  \uses{def:physics:phyx-mini-0752:phyxminiproblems-problemphyxmini0752-answerchoice}
  Answer label recorded in the supplied dataset metadata
\end{definition}
\begin{proof}
  This is the declared model definition of recorded Answer Choice.
\end{proof}
\begin{definition}[answer Choice Error Meters Per Second]
  \label{def:physics:phyx-mini-0752:phyxminiproblems-problemphyxmini0752-answerchoiceerrormeterspersecond}
  \lean{PhyXMiniProblems.ProblemPhyXMini0752.answerChoiceErrorMetersPerSecond}
  \uses{def:physics:phyx-mini-0752:phyxminiproblems-problemphyxmini0752-speedinmeterspersecond, def:physics:phyx-mini-0752:phyxminiproblems-problemphyxmini0752-volcanicbombsetup, def:physics:phyx-mini-0752:phyxminiproblems-problemphyxmini0752-answerchoice}
  Absolute display error for a candidate answer speed
\end{definition}
\begin{proof}
  This is the declared model definition of answer Choice Error Meters Per Second.
\end{proof}
\begin{definition}[Matches Displayed Launch Speed]
  \label{def:physics:phyx-mini-0752:phyxminiproblems-problemphyxmini0752-matchesdisplayedlaunchspeed}
  \lean{PhyXMiniProblems.ProblemPhyXMini0752.MatchesDisplayedLaunchSpeed}
  \uses{def:physics:phyx-mini-0752:phyxminiproblems-problemphyxmini0752-volcanicbombsetup, def:physics:phyx-mini-0752:phyxminiproblems-problemphyxmini0752-answerchoice, def:physics:phyx-mini-0752:phyxminiproblems-problemphyxmini0752-answerchoiceerrormeterspersecond}
  Agreement with a speed displayed to the nearest tenth of a metre per second
\end{definition}
\begin{proof}
  This is the declared model definition of Matches Displayed Launch Speed.
\end{proof}
\begin{definition}[Is Unique Closest Launch Speed Choice]
  \label{def:physics:phyx-mini-0752:phyxminiproblems-problemphyxmini0752-isuniqueclosestlaunchspeedchoice}
  \lean{PhyXMiniProblems.ProblemPhyXMini0752.IsUniqueClosestLaunchSpeedChoice}
  \uses{def:physics:phyx-mini-0752:phyxminiproblems-problemphyxmini0752-volcanicbombsetup, def:physics:phyx-mini-0752:phyxminiproblems-problemphyxmini0752-answerchoice, def:physics:phyx-mini-0752:phyxminiproblems-problemphyxmini0752-answerchoiceerrormeterspersecond}
  The selected displayed speed is strictly closer than every other choice
\end{definition}
\begin{proof}
  This is the declared model definition of Is Unique Closest Launch Speed Choice.
\end{proof}
\begin{lemma}[launch Speed eq required Expression]
  \label{lem:physics:phyx-mini-0752:phyxminiproblems-problemphyxmini0752-launchspeed-eq-requiredexpression}
  \lean{PhyXMiniProblems.ProblemPhyXMini0752.launchSpeed_eq_requiredExpression}
  \uses{def:physics:phyx-mini-0752:phyxminiproblems-problemphyxmini0752-speedinmeterspersecond, def:physics:phyx-mini-0752:phyxminiproblems-problemphyxmini0752-volcanicbombsetup, def:physics:phyx-mini-0752:phyxminiproblems-problemphyxmini0752-matchesproblemdescription, def:physics:phyx-mini-0752:phyxminiproblems-problemphyxmini0752-usesstandardnearearthgravity, def:physics:phyx-mini-0752:phyxminiproblems-problemphyxmini0752-hasphysicalprojectileparameters, def:physics:phyx-mini-0752:phyxminiproblems-problemphyxmini0752-satisfiesuniformgravityprojectilelaws, def:physics:phyx-mini-0752:phyxminiproblems-problemphyxmini0752-requiredlaunchspeedinmeterspersecond}
  The governing equations determine the initial speed's exact unrounded SI readout. This conclusion is not a field of any premise structure
\end{lemma}
\begin{proof}
  The conclusion follows from the governing relations and figure readouts named in the statement of launch Speed eq required Expression.
\end{proof}
% --- Archon declaration topology end ---
% --- Archon physics formalization source end ---
